% !TeX program = lualatex
% Compile twice: lualatex 170.tex
% All references and prose are embedded; no BibTeX database or external inputs.
\documentclass[11pt,a4paper]{article}
\usepackage[margin=24mm,headheight=14pt]{geometry}
\usepackage{fontspec}
\setmainfont{Latin Modern Roman}
\setsansfont{Latin Modern Sans}
\setmonofont{Latin Modern Mono}
\usepackage{amsmath,amssymb,mathtools}
\usepackage{unicode-math}
\setmathfont{Latin Modern Math}
\usepackage{microtype}
\usepackage{enumitem,xurl,needspace}
\usepackage[dvipsnames]{xcolor}
\usepackage{hyperref}
\hypersetup{unicode=true,colorlinks=true,linkcolor=MidnightBlue,urlcolor=MidnightBlue,
 pdftitle={500 Mathematical Problem Statements},pdfauthor={Source-annotated compilation},bookmarksnumbered=true}
\usepackage{bookmark}
\usepackage{tocloft}
\setlength{\cftsecnumwidth}{3em}
\cftsetpnumwidth{2.5em}
\cftsetrmarg{4em}
\usepackage{titlesec}
\titleformat{\section}{\Large\bfseries\raggedright}{\thesection}{0.7em}{}
\usepackage{fancyhdr}
\pagestyle{fancy}\fancyhf{}
\fancyhead[L]{\small\sffamily Mathematical problem statements}
\fancyhead[R]{\small Source edition 22}
\fancyfoot[C]{\thepage}
\setlength{\parindent}{0pt}
\setlength{\parskip}{5pt plus 1pt}
\setlength{\emergencystretch}{3em}
\setcounter{tocdepth}{1}
\setcounter{secnumdepth}{1}
\setlist[itemize]{leftmargin=3em,itemsep=5pt,topsep=4pt,parsep=0pt}
\allowdisplaybreaks
\newcommand{\N}{\mathbb{N}}
\newcommand{\Z}{\mathbb{Z}}
\newcommand{\Q}{\mathbb{Q}}
\newcommand{\R}{\mathbb{R}}
\newcommand{\C}{\mathbb{C}}
\newcommand{\F}{\mathbb{F}}
\newcommand{\A}{\mathbb{A}}
\newcommand{\PP}{\mathbb P}
\newcommand{\E}{\mathbb{E}}
\newcommand{\eps}{\varepsilon}
\newcommand{\dd}{\,\mathrm d}
\newcommand{\sref}[2]{\hyperlink{source:#1:#2}{[S#2]}}
\newcommand{\eref}[2]{\hyperlink{extra:#1:#2}{[E#2]}}
% Turn the next switch off for a shorter reading copy. The quotations preserve
% every exactTarget field, including qualifications omitted from a short summary.
\newif\ifshowcatalogue
\showcataloguetrue
\newcommand{\cataloguescope}[1]{\ifshowcatalogue
 \paragraph{Exact scope in the supplied catalogue.}
 \begin{quote}\small #1\end{quote}\fi}

% Standalone notebook wrapper; original problem section retained verbatim.
\numberwithin{equation}{section}
\hypersetup{pdftitle={Research attempt -- problem 170},
 pdfauthor={Research notebook; source compilation retained}}
\fancyhead[L]{\small\sffamily Research notebook}
\fancyhead[R]{\small\sffamily Problem 170 | 27 September 2026}
\begin{document}
\setcounter{section}{169}
% ================================================================
% problemId: problem.dirichlet-divisor-problem
% source releaseRank: 170; source releaseStatus: open_disputed_claim
% exactTarget SHA-256: e100bd8a6f8495f8a8b1d88f0db13236b9502e26cc2ad23c836cbab1edb218d9
\clearpage
\section{\texorpdfstring{Dirichlet Divisor Problem}{Dirichlet Divisor Problem}}\label{170}
\subsection{Definitions and mathematical statement}
For real $x\ge2$, let $d(n)=\sum_{a\mid n}1$ and
\[
 \Delta(x)=\sum_{1\le n\le x}d(n)-x\log x-(2\gamma-1)x,
 \qquad \gamma=\lim_{N\to\infty}\left(\sum_{n=1}^N\frac1n-\log N\right).
\]
The exact selected problem is to determine
$\alpha_2=\inf\{a\in\R:\forall\eps>0,\ \Delta(x)=O_\eps(x^{a+\eps})\}$.
Thus the implied constant may depend on $\eps$ but not on $x$. Specifying this infimum is not the same as proving a bound with no $\eps$ loss at its endpoint, nor is a mean-square bound a pointwise bound. The catalogue asks for the optimal exponent rather than explicitly adding an endpoint estimate.
\par\smallskip\textit{Formulation and concept references:} \sref{170}{1}.
\cataloguescope{Let d(n) count the positive divisors of n and Delta(x)=sum\_\{n<=x\} d(n)-x log x-(2 gamma-1)x for real x>=2. Determine the optimal exponent alpha\_2=inf\{a: for every epsilon>0, Delta(x)=O\_epsilon(x\textasciicircum{}(a+epsilon)) as x tends to infinity\}.}
\subsection{Short English statement}
How small can the exponent in a uniform pointwise error estimate for the summatory divisor function be?
% BEGIN RESEARCH ADDITION ATTEMPT-170
\subsection{Research attempt}

\paragraph{Current frontier.}
The lower-bound scale forces $\alpha_2\ge1/4$. The introduction to
\sref{170}{1} records the upper bound $0.314483\ldots$ of Li--Yang and the
stronger-than-$x^{1/4}$ logarithmic omega results; neither determines the
infimum. The additional catalogue references include a smoothed problem
\sref{170}{2}, a paper claiming a proof \sref{170}{3}, an optimal-shifting
manuscript \sref{170}{4}, and an omega-result manuscript \sref{170}{5}.
They have been inspected as status sources, not accepted as proof
certificates. A smoothed upper bound or an omega estimate is not the missing
pointwise upper bound. The analysis below uses no claimed resolution from
those papers.

\paragraph{Attack route.}
A Voronoi-sum route needs cancellation in oscillatory sums uniformly at each
$x$; a smoothing route must pay for removal of the averaging. A high-moment
route can exploit the monotonicity of the summatory function to detect a
large pointwise value on a short interval. We pursue the latter two routes
together. The aim is an exact accounting of the power lost in converting
average control into pointwise control, rather than an assertion that
square-root cancellation is automatic.

\paragraph{Attempt: a large value must persist on one side.}
Write $D(t)=\sum_{n\le t}d(n)$ and
$F(t)=t\log t+(2\gamma-1)t$. On $[X/2,3X]$, for $X$ sufficiently large,
$0<F'(t)=\log t+2\gamma\le L_X$, where $L_X=C\log X$ for an absolute
$C$. The elementary hyperbola identity
\[
 D(t)=2\sum_{n\le\sqrt t}\lfloor t/n\rfloor-\lfloor\sqrt t\rfloor^2
\]
and the harmonic-sum estimate give $\Delta(t)=O(\sqrt t)$.
To check the latter, replace each floor by $t/n+O(1)$ and use
$\sum_{n\le m}n^{-1}=\log m+\gamma+O(m^{-1})$ with
$m=\lfloor\sqrt t\rfloor$; expanding $m^2=t+O(\sqrt t)$ supplies the
stated main term and error. This crude bound is used only to keep the next
short interval inside $[X/2,3X]$.

Fix $x\in[X,2X]$, and let $V=|\Delta(x)|>0$ and $h=V/(2L_X)$.
If $\Delta(x)=V$, monotonicity of $D$ gives, for $x\le t\le x+h$,
\[
 \Delta(t)\ge\Delta(x)-\bigl(F(t)-F(x)\bigr)\ge V/2.
\]
If $\Delta(x)=-V$, use $x-h\le t\le x$ instead, obtaining
$\Delta(t)\le-V/2$. This works at integer $x$ as well; it uses the actual
right-continuous counting function and no differentiability of $D$.
For every fixed real $q>0$ we have therefore proved
\begin{equation}\label{eq170-persistence}
 \int_{X/2}^{3X}|\Delta(t)|^q\,dt
       \ge \frac{V^{q+1}}{2^{q+1}L_X}.
\end{equation}
The interval may be exceedingly short. Its guaranteed length is $V/\log X$,
not $X$ or $\sqrt X$.

\paragraph{Attempt: the exact high-moment reduction.}
Consider the following family of estimates, explicitly left conditional:
\begin{equation}\label{eq170-moment}
 \int_{X/2}^{3X}|\Delta(t)|^q\,dt
       \ll_{q,\eta}X^{1+q/4+\eta}
       \quad\text{for every fixed }q>0,\ \eta>0.
\end{equation}
Combining it with \eqref{eq170-persistence} gives
\begin{equation}\label{eq170-power}
 \sup_{X\le x\le2X}|\Delta(x)|
 \ll_{q,\eta}(\log X)^{1/(q+1)}
 X^{\frac14+\frac{3}{4(q+1)}+\frac{\eta}{q+1}}.
\end{equation}
For any prescribed $\eps>0$, first choose a fixed $q$ so large that
$3/(4(q+1))<\eps/3$, then choose $\eta$ and absorb the logarithm.
A dyadic covering yields $\Delta(x)=O_\eps(x^{1/4+\eps})$.
Together with the known lower bound this determines $\alpha_2=1/4$.
There is no need for constants uniform as $q\to\infty$: the choice of $q$
is made before taking $X\to\infty$.

Conversely, the collection of pointwise bounds
$\Delta(t)=O_\eps(t^{1/4+\eps})$ gives \eqref{eq170-moment} by taking
$\eps=\eta/q$ and integrating over an interval of length $O(X)$.
Thus \eqref{eq170-moment} for arbitrarily large fixed $q$ is equivalent to
the expected value of the exponent, in the presence of the recorded lower
bound. It is not a strictly easier theorem that has now been proved.
The elementary deduction exposes the quantitative input a high-moment
strategy must deliver.

For comparison, at $q=2$ the exponent in \eqref{eq170-power} is $1/2$,
not $1/4$. Even an optimal second moment does not give the pointwise target
through this argument. Values with height $V$ and persistence length only
$V/\log X$ cost about $V^3/\log X$ in that moment. The integral can permit
rare large values while retaining its conjectured scale.

\paragraph{Attempt: unsmoothing with a controlled loss.}
For $0<h\le X/4$, define one-sided means
\[
 A_+(x,h)=h^{-1}\int_x^{x+h}\Delta(t)\,dt,\qquad
 A_-(x,h)=h^{-1}\int_{x-h}^{x}\Delta(t)\,dt.
\]
The same monotonicity argument, integrated before any estimate is applied,
gives the exact inequalities
\begin{equation}\label{eq170-unsmooth}
 A_-(x,h)-L_Xh/2\le\Delta(x)
                    \le A_+(x,h)+L_Xh/2.
\end{equation}
In particular, a bound $|A_\pm(x,h)|\ll X^{1/4+\eps}$ is useful at the
expected pointwise scale only if the chosen width also satisfies
$h\log X\ll X^{1/4+\eps}$. Obtaining such a bound only at widths of
order $X^{1/2}$ would leave an error of order $X^{1/2}\log X$.
An arbitrary signed smoothing kernel need not support either comparison;
its mass, positivity, support, and first moment must be tracked separately.
This is why a smoothed statement, such as the distinct problem named in
\sref{170}{2}, cannot be inserted directly as an answer to the catalogue.

One possible next input is a two-scale estimate proving the short averages
in \eqref{eq170-unsmooth} from a longer smoothing and a uniform increment
bound. However, the needed increment bound would itself have to control
short divisor sums after subtraction of their main term. No such estimate
at the required scale is established in this attempt. The elementary bound
on the number of terms would recover only the loss already displayed.

\paragraph{Stress test.}
The signs in the one-sided argument are essential: positive excursions
persist to the right, negative excursions to the left. Attempting a
symmetric persistence statement would incorrectly control the jumps of
$D$. The width and logarithmic loss have not been discarded, and no
interchange of an infinite sum and integral is used.

Finally, determining an infimum is not the same as proving a literal
endpoint bound $O(x^{1/4})$. The lower results cited in \sref{170}{1}
already include growing logarithmic factors along suitable sequences. The
conditional argument promises only every $\eps$-loss bound, exactly as
required for the infimum. It supplies neither an endpoint estimate nor a
new omega theorem. None of the recent proof claims is independently
validated by these elementary comparisons.

\paragraph{Outcome.}
\textbf{Outcome: Conditional reduction.}

The attempt proves an explicit moment-to-supremum inequality and exact
one-sided unsmoothing inequalities. It identifies arbitrarily high fixed
moments at the expected scale as a sufficient and, here, equivalent missing
input for $\alpha_2=1/4$. It also quantifies why second moments and long
smoothings do not close the argument. These are deductions and diagnostic
lemmas, not a new bound on $\alpha_2$ or a claim of literature priority.
% END RESEARCH ADDITION ATTEMPT-170
\subsection{Sources}
\begin{itemize}\raggedright
\item[\textnormal{[S1]}]\hypertarget{source:170:1}{} Karak and Mahatab, The Piltz divisor Problem in Number Fields Using The Resonance Method (2025; journal 2026).
\url{https://arxiv.org/html/2506.22587v1}.
{\small\textit{Catalogue formulation source.}}
\item[\textnormal{[S2]}]\hypertarget{source:170:2}{} Status source for Dirichlet Divisor Problem.
\url{https://arxiv.org/html/2601.01905v2}.
{\small\textit{Catalogue research/status reference; not a proof endorsement.}}
\item[\textnormal{[S3]}]\hypertarget{source:170:3}{} Status source for Dirichlet Divisor Problem.
\url{https://arxiv.org/pdf/1105.6155}.
{\small\textit{Catalogue research/status reference; not a proof endorsement.}}
\item[\textnormal{[S4]}]\hypertarget{source:170:4}{} Status source for Dirichlet Divisor Problem.
\url{https://arxiv.org/html/2604.18624v2}.
{\small\textit{Catalogue research/status reference; not a proof endorsement.}}
\item[\textnormal{[S5]}]\hypertarget{source:170:5}{} Status source for Dirichlet Divisor Problem.
\url{https://arxiv.org/html/2605.21476v1}.
{\small\textit{Catalogue research/status reference; not a proof endorsement.}}
\end{itemize}

\end{document}
